\documentclass{article} % Starts a article
\usepackage{amsmath} % Imports amsmath
\title{Aula 2} % Title
\date{2020 -08 -26}
\author{Lucas Tatsuya Tanaka}
\begin{document} % Begins a document
  \maketitle

    \section{Qual a importancia do Software}
    \begin{itemize}
	\item Levando em consideracao um mundo globalizado e muito conectado 
            a informacao se tornou um novo comodite
	\item Uma mudanca princiapal de comutacao por circuitos para comunicacao por pacotes 
            ou seja mudanca de circuitos para pacotes como base em internet 
	\item Com o surgimento da internet antigas armas de conhecimento como a criptogerafia
            se tornaram em ferramentas para o uso cotidiano na sociedade 
    \end{itemize}

    \section{Riscos do Softaware}
    \begin{itemize}
        \item o tamanho implica em Risco quanto maior ele deve maiores sao as chances de 
            serem abandonados caso sejam pequenas podem ser modificados
    \end{itemize}

    \section{Dificuldade de fazer Software}
    \begin{itemize}
        \item com o oumento dos computadores maiores os problemas se tornaram ocorreu um 
            aumento da dificuldade dos problemas
        \item A lei de Brooks-Adicionar pessoas a um projeto atrasado atrasa mais ainda o 
            projeto
    \end{itemize}

    \section{Referencias}
    Edsger W. Dijkstra (Ganhador do premio Turing)- Humble programer 
    Fred Brooks- The Mythical Man-Month (Criador da Lei de Brooks)

numeros de ponto de funcao
\end{document}
